\documentclass[answers,12pt,addpoints]{exam} 
\usepackage{import}

\import{C:/Users/prana/OneDrive/Desktop/MathNotes}{style.tex}

% Header
\newcommand{\name}{Pranav Tikkawar}
\newcommand{\course}{Experimental Design}
\newcommand{\assignment}{Exam 1 (Take-home)}
\author{\name}
\title{\course \ - \assignment}

\begin{document}
\maketitle

\textbf{Problem 1}

In a study of starting salaries for assistant professors, five male
assistant professors at each of three types of institutions were randomly
polled and their starting salaries (in \$1000) were recorded. The summary
statistics are:

\[
\begin{array}{lccc}
\hline
 & \text{Public} & \text{Private} & \text{Church-affiliated} \\
\hline
\text{Mean}     & 53.4 & 70.8 & 56.1 \\
\text{Variance} & 72.8 & 48.2 & 60.3 \\
\hline
\end{array}
\]

There are $n_i = 5$ observations per group ($a=3$, $N=15$).

\medskip
\textbf{(a) Type of experimental design and model}

This is a single-factor completely randomized design (CRD) with one factor
``institution type'' at three levels. A fixed-effects one-way ANOVA model is
\[
Y_{ij} = \mu + \tau_i + \varepsilon_{ij},\quad i=1,2,3,\; j=1,\dots,5,
\]
where $\mu$ is the overall mean, $\tau_i$ is the effect of institution type $i$
with the usual constraint $\sum_{i=1}^3 \tau_i = 0$, and
$\varepsilon_{ij} \sim \mathrm{NID}(0,\sigma^2)$.

\medskip
\textbf{(b) Test for differences in average starting salaries}

We test
\[
H_0:\,\mu_1 = \mu_2 = \mu_3
\quad\text{vs}\quad
H_a:\,\text{at least one }\mu_i\text{ differs}.
\]

The grand mean is
\[
\bar{y}_{..} = \frac{53.4 + 70.8 + 56.1}{3} = 60.1.
\]

The treatment sum of squares is
\[
SS_{\text{Trt}}
= n\sum_{i=1}^3(\bar{y}_i - \bar{y}_{..})^2
= 5\bigl[(53.4-60.1)^2 + (70.8-60.1)^2 + (56.1-60.1)^2\bigr]
= 876.9.
\]

The error sum of squares is
\[
SS_E = \sum_{i=1}^3 (n_i - 1)s_i^2
= 4(72.8 + 48.2 + 60.3) = 725.2,
\]
so the total sum of squares is
\[
SS_T = SS_{\text{Trt}} + SS_E = 876.9 + 725.2 = 1602.1.
\]

Degrees of freedom and mean squares:
\[
df_{\text{Trt}} = a-1 = 2,\quad
df_E = N-a = 12,\quad
df_T = N-1 = 14,
\]
\[
MS_{\text{Trt}} = \frac{SS_{\text{Trt}}}{2} = 438.45,\qquad
MS_E = \frac{SS_E}{12} = 60.4333.
\]

The ANOVA $F$-statistic is
\[
F = \frac{MS_{\text{Trt}}}{MS_E}
= \frac{438.45}{60.4333} \approx 7.26,
\]
which is compared to an $F$ distribution with $(2,12)$ degrees of freedom.
The corresponding $p$-value from R is
\[
p = 0.0086.
\]
Since $p < 0.05$, we reject $H_0$ and conclude that there are significant
differences in average starting salaries among the three institution types.

\medskip
\textbf{(c) Estimate of the standard deviation of the error}

The error mean square $MS_E$ is the usual unbiased estimator of $\sigma^2$.
Hence
\[
\hat{\sigma}^2 = MS_E = 60.4333,
\qquad
\hat{\sigma} = \sqrt{MS_E} \approx 7.77,
\]
so the estimated error standard deviation is about \$7{,}770.

\medskip
\textbf{(d) Pairwise comparisons using Tukey HSD and Scheff\'e}

With a balanced design and common replication $n=5$, the standard error for
pairwise differences of treatment means is
\[
SE(\bar{y}_i - \bar{y}_j) =
\sqrt{\frac{2 MS_E}{n}}
= \sqrt{\frac{2 \times 60.4333}{5}}
= \sqrt{24.1733} \approx 4.92.
\]

The pairwise sample mean differences (treatment $j$ minus treatment $i$) are
\[
\bar{y}_{\text{Priv}} - \bar{y}_{\text{Pub}} = 17.4,\quad
\bar{y}_{\text{Ch}}   - \bar{y}_{\text{Pub}} = 2.7,\quad
\bar{y}_{\text{Ch}}   - \bar{y}_{\text{Priv}} = -14.7.
\]

\emph{Tukey HSD.}
For $\alpha = 0.05$, $a = 3$, and $df_E = 12$, the studentized range critical
value from R is
\[
q_{0.05,3,12} = 3.7729,
\]
so
\[
HSD = q_{0.05,3,12}\sqrt{\frac{MS_E}{n}}
= 3.7729 \sqrt{\frac{60.4333}{5}}
\approx 3.7729 \times 3.48 \approx 13.1.
\]
A pair $(i,j)$ is significantly different if $|\bar{y}_i - \bar{y}_j| > HSD$.

\begin{itemize}
  \item Private vs.\ Public:
  $|\bar{y}_{\text{Priv}} - \bar{y}_{\text{Pub}}|
  = 17.4 > 13.1$ $\Rightarrow$ significant.
  \item Church vs.\ Public:
  $|\bar{y}_{\text{Ch}} - \bar{y}_{\text{Pub}}|
  = 2.7 < 13.1$ $\Rightarrow$ not significant.
  \item Church vs.\ Private:
  $|\bar{y}_{\text{Ch}} - \bar{y}_{\text{Priv}}|
  = 14.7 > 13.1$ $\Rightarrow$ significant.
\end{itemize}

Thus, under Tukey's method, Private institutions have significantly higher
starting salaries than both Public and Church-affiliated institutions, while
Public and Church-affiliated institutions do not differ significantly.

\emph{Scheff\'e method.}
For a one-way ANOVA with $a=3$ groups, Scheff\'e's critical factor at
$\alpha = 0.05$ is
\[
k_S = (a-1)F_{0.05,\,a-1,\,df_E}
= 2 F_{0.05,2,12}.
\]
From R,
\[
F_{0.05,2,12} = 3.8853,\quad
k_S = 2 \times 3.8853 \approx 7.77.
\]

For a pairwise contrast $C = \mu_i - \mu_j$, the estimated contrast is
$\hat{C} = \bar{y}_i - \bar{y}_j$ with
\[
SE(\hat{C}) = SE(\bar{y}_i - \bar{y}_j) \approx 4.92,
\]
and Scheff\'e's test rejects $H_0: C=0$ if
\[
\frac{\hat{C}^2}{SE(\hat{C})^2} > k_S.
\]

Using the R output,
\[
\begin{aligned}
T^2_{\text{Priv-Pub}} &= 12.52,\\
T^2_{\text{Ch-Pub}}   &= 0.30,\\
T^2_{\text{Ch-Priv}}  &= 8.94,
\end{aligned}
\]
so
\[
T^2_{\text{Priv-Pub}} > 7.77,\quad
T^2_{\text{Ch-Priv}} > 7.77,\quad
T^2_{\text{Ch-Pub}} < 7.77.
\]
Therefore, Scheff\'e's method also finds that Private vs.\ Public and
Private vs.\ Church-affiliated are significantly different, while Public vs.\
Church-affiliated is not. In this example, Tukey and Scheff\'e give the same
pairwise conclusions.

\medskip
\textbf{(e) Kruskal--Wallis test}

A Kruskal--Wallis test with $a=3$ groups has $df = a-1 = 2$ and, under $H_0$,
the test statistic is approximately $\chi^2_2$-distributed. We are told that
the Kruskal--Wallis test statistic is $H = 8.09$. The 5\% upper critical value
for $\chi^2_2$ is about $5.99$, and since $8.09 > 5.99$, the Kruskal--Wallis
test also rejects $H_0$ at the 5\% level. Thus the nonparametric analysis
confirms the ANOVA result: there is evidence of differences in starting salaries
among the three institution types, so our conclusion in part (b) does not
change.

\medskip
\textbf{(f) Possible improved model}

The current model treats only institution type as a factor and ignores other
important sources of variability. An improved model might:

\begin{itemize}
  \item include additional factors or covariates such as academic discipline,
  geographic region, or research vs.\ teaching focus, which are known to affect
  starting salaries;
  \item use blocking on departments, regions, or year of hire to control for
  nuisance variation via a randomized complete block design;
  \item consider transformations (e.g., modeling log-salary) if residual plots
  indicate skewness or nonconstant variance across institution types.
\end{itemize}

Such extensions would likely reduce unexplained variability and provide a more
realistic assessment of the effect of institution type on starting salaries.

\textbf{Problem 2}

Three groups ($\#1,\#2,\#3$) of students were subjected to different teaching
techniques and tested at the end of the course. There are $n_i = 6$ students in
each group, and the average scores are
\[
\bar{y}_1 = 72.3,\quad \bar{y}_2 = 79.5,\quad \bar{y}_3 = 65.8.
\]
The overall variance of all 18 scores is 84.03, and we are told that the error
variance estimate from the ANOVA model is
\[
MS_E = 57.84.
\]

\medskip
\textbf{(a) Complete ANOVA table}

The grand mean is
\[
\bar{y}_{..} = \frac{72.3 + 79.5 + 65.8}{3} = 72.5333.
\]

The treatment sum of squares is
\[
SS_{\text{Trt}}
= n\sum_{i=1}^3(\bar{y}_i - \bar{y}_{..})^2
= 6\bigl[(72.3-72.5333)^2 + (79.5-72.5333)^2 + (65.8-72.5333)^2\bigr]
= 563.56.
\]

We are given $MS_E = 57.84$. The error degrees of freedom are
\[
df_E = N-a = 18-3 = 15,
\]
so
\[
SS_E = MS_E \times df_E = 57.84 \times 15 = 867.6.
\]
Thus the total sum of squares is
\[
SS_T = SS_{\text{Trt}} + SS_E = 563.56 + 867.6 = 1431.16.
\]

Degrees of freedom and mean squares:
\[
df_{\text{Trt}} = a-1 = 2,\quad df_E = 15,\quad df_T = N-1 = 17,
\]
\[
MS_{\text{Trt}} = \frac{SS_{\text{Trt}}}{2} = 281.78,\qquad
MS_E = 57.84.
\]

The ANOVA $F$-statistic is
\[
F = \frac{MS_{\text{Trt}}}{MS_E}
= \frac{281.78}{57.84} \approx 4.87.
\]

The ANOVA table is therefore

\[
\begin{array}{lrrrr}
\hline
\text{Source} & \text{df} & \text{SS} & \text{MS} & F \\
\hline
\text{Treatment} & 2  & 563.56 & 281.78 & 4.87 \\
\text{Error}     & 15 & 867.60 & 57.84  &      \\
\text{Total}     & 17 & 1431.16&        &      \\
\hline
\end{array}
\]

\medskip
\textbf{(b) Conclusion about differences in teaching techniques}

We test
\[
H_0:\,\mu_1 = \mu_2 = \mu_3
\quad\text{vs}\quad
H_a:\,\text{at least one }\mu_i\text{ differs}.
\]

From the ANOVA table,
\[
F = 4.8717,\quad df_1 = 2,\quad df_2 = 15.
\]
The corresponding $p$-value from R is
\[
p = 0.0234.
\]
Since $p < 0.05$, we reject $H_0$ at the 5\% level and conclude that there are
significant differences in mean test scores among the three teaching
techniques.

\medskip
\textbf{(c) Comparison of treatments to the control (group 2)}

Suppose group 2 is the control and groups 1 and 3 are treatments. The
differences in sample means relative to the control are
\[
\bar{y}_1 - \bar{y}_2 = 72.3 - 79.5 = -7.2,\qquad
\bar{y}_3 - \bar{y}_2 = 65.8 - 79.5 = -13.7.
\]

The standard error for the difference between a treatment mean and the control
mean (with equal group sizes $n=6$) is
\[
SE_{\text{diff}} = \sqrt{\frac{2 MS_E}{n}}
= \sqrt{\frac{2 \times 57.84}{6}} = 4.3909.
\]

Thus the $t$-statistics (with $df = 15$) are
\[
t_{1-2} = \frac{-7.2}{4.3909} = -1.64,\qquad
t_{3-2} = \frac{-13.7}{4.3909} = -3.12.
\]

At the usual $\alpha = 0.05$ level for a two-sided test, the comparison of
group 3 to group 2 is significant (group 3 has a significantly lower mean
score than the control), while the comparison of group 1 to group 2 is not
statistically significant, even though group~1's mean is somewhat lower than
the control.

\medskip
\textbf{(d) Average treatment effect vs control with Bonferroni correction}

We now test whether the \emph{average} of the two treatment groups differs from
the control. Consider the contrast
\[
C = \tfrac{1}{2}\mu_1 + \tfrac{1}{2}\mu_3 - \mu_2.
\]
Its estimate is
\[
\hat{C} = \tfrac{1}{2}(72.3 + 65.8) - 79.5
= \tfrac{1}{2}(138.1) - 79.5 = 69.05 - 79.5 = -10.45.
\]

For equal group sizes $n=6$, the variance of $\hat{C}$ is
\[
\operatorname{Var}(\hat{C}) =
MS_E\left(\frac{0.5^2}{n} + \frac{0.5^2}{n} + \frac{1^2}{n}\right)
= MS_E\cdot\frac{1.5}{n},
\]
so
\[
SE_C = \sqrt{\frac{1.5\,MS_E}{n}}
= \sqrt{\frac{1.5 \times 57.84}{6}} = 3.8026.
\]

The corresponding $t$-statistic (with $df = 15$) is
\[
t = \frac{\hat{C}}{SE_C} = \frac{-10.45}{3.8026} = -2.75.
\]

We are asked to test this contrast using a Bonferroni correction assuming a
total of $K = 10$ post-hoc tests. With overall $\alpha = 0.05$, the
Bonferroni-adjusted per-test level is
\[
\alpha^* = \frac{\alpha}{K} = \frac{0.05}{10} = 0.005.
\]
For a two-sided test, we use $\alpha^*/2 = 0.0025$ in each tail. The critical
$t$ value from R is
\[
t_{1-\alpha^*/2,15} = t_{0.9975,15} = 3.286.
\]

Since
\[
|t| = 2.75 < 3.286,
\]
we fail to reject $H_0 : C = 0$ at the Bonferroni-adjusted level. Thus, when
controlling for a family of 10 post-hoc tests, the average of the two
treatment groups is not significantly different from the control group, even
though the unadjusted $t$-statistic suggests a noticeable decrease in score.

\textbf{Problem 3 \,(DAE 4.17)}

The response is gene expression measured for three leukemia treatments:
MP only, MP + HDMTX, and MP + LDMTX, with $n=10$ observations per
treatment ($N=30$). The data are in \texttt{Mid1Prob3.csv}.

\medskip
\textbf{(a) Test whether treatment means differ (raw data)}

A one-way fixed-effects ANOVA model is
\[
Y_{ij} = \mu + \tau_i + \varepsilon_{ij},\quad
i=1,2,3,\; j=1,\dots,10,\qquad
\varepsilon_{ij} \sim \mathrm{NID}(0,\sigma^2),
\]
where $i$ indexes treatment.

Sample means and standard deviations by treatment are
\[
\begin{array}{lccc}
\hline
\text{Treatment} & n & \text{Mean} & \text{SD} \\
\hline
\text{MP + HDMTX} & 10 & 479.3 & 371.0 \\
\text{MP + LDMTX} & 10 & 165.0 & 126.0 \\
\text{MP Only}    & 10 & 238.0 & 308.0 \\
\hline
\end{array}
\]

From the ANOVA on the raw responses (R command
\texttt{aov(Gene.Expression $\sim$ Treatment)}),
\[
\begin{array}{lrrrr}
\hline
\text{Source} & \text{Df} & \text{SS} & \text{MS} & F \\
\hline
\text{Treatment} &  2 & 538{,}442 & 269{,}221 & 3.25 \\
\text{Error}     & 27 & 2{,}236{,}974 & 82{,}851 &     \\
\hline
\end{array}
\]
with $p$–value $p = 0.0544$.

At $\alpha = 0.05$, $p$ is slightly larger than 0.05, so we \emph{do not}
reject $H_0$ at the 5\% level based strictly on this ANOVA. The evidence
for a difference in mean gene expression across treatments is marginal
but not conventionally significant.

\medskip
\textbf{(b) Normality check for raw data}

The residuals from the raw-data ANOVA show substantial deviation from
normality. The Shapiro–Wilk test applied to the raw residuals gives
\[
W = 0.9165,\quad p = 0.0218,
\]
which is below 0.05, indicating a statistically significant departure
from normality. The residual histogram and QQ plot 
exhibit pronounced right skew and a few large positive residuals,
consistent with the presence of outliers or heavy right tails. Thus the
normality assumption is questionable on the original scale.

\begin{center}
\includegraphics[width=0.45\textwidth]{img/3b.png}
\end{center}

\medskip
\textbf{(c) Log transform and ANOVA on transformed data}

Because the raw data are positive and strongly right-skewed, a log
transform is natural. Let
\[
Z_{ij} = \log(Y_{ij}),
\]
and fit the one-way ANOVA model
\[
Z_{ij} = \mu^\ast + \tau_i^\ast + \varepsilon_{ij}^\ast.
\]

On the log scale, the sample means and standard deviations are
\[
\begin{array}{lccc}
\hline
\text{Treatment} & n & \text{Mean} & \text{SD} \\
\hline
\text{MP + HDMTX} & 10 & 5.73 & 1.13 \\
\text{MP + LDMTX} & 10 & 4.74 & 1.00 \\
\text{MP Only}    & 10 & 4.79 & 1.18 \\
\hline
\end{array}
\]

The ANOVA on the log-transformed responses (R command
\texttt{aov(logGE $\sim$ Treatment)}) yields
\[
\begin{array}{lrrrr}
\hline
\text{Source} & \text{Df} & \text{SS} & \text{MS} & F \\
\hline
\text{Treatment} &  2 &  6.297 & 3.1485 & 2.57 \\
\text{Error}     & 27 & 33.069 & 1.2248 &      \\
\hline
\end{array}
\]
with $p = 0.0951$.

Thus, on the log scale the evidence for treatment differences is weaker:
at $\alpha = 0.05$ we again fail to reject $H_0$; the $p$–value is
between 0.05 and 0.10, so one might describe this as only marginal
evidence of differences in mean log gene expression.

\medskip
\textbf{(d) Residual analysis for the transformed model}

For the log-scale ANOVA, we examine the residuals. The residuals vs
fitted plot for the log model shows a fairly symmetric cloud of points
with no strong funnel shape; the variability appears roughly constant
across fitted values. The QQ plot of the log-model residuals lies much
closer to a straight line than for the raw data.

The Shapiro–Wilk test on the log-model residuals gives
\[
W = 0.9565,\quad p = 0.2511,
\]
which is not significant; we do not reject normality at the 5\% level.
This suggests that the log transform greatly improves the normality of
the errors and makes the ANOVA assumptions more plausible.

In summary, the log transformation yields a model with much better
residual behavior (approximately normal and homoscedastic), but even on
this scale the treatment effect is not statistically significant at the
$\alpha = 0.05$ level.

\begin{center}
\includegraphics[width=0.45\textwidth]{img/3d.png}
\end{center}

\newpage
\textbf{R Code}
\begin{verbatim}
> ########################
> # Problem 1
> ########################
>
> # 1(a-d): One-way ANOVA from summary stats
> means  <- c(Public = 53.4, Private = 70.8, Church = 56.1)
> vars   <- c(Public = 72.8, Private = 48.2, Church = 60.3)
> n      <- 5
> grand_mean <- mean(means)
> grand_mean
[1] 60.1
> SS_trt <- n * sum((means - grand_mean)^2)
> SS_trt
[1] 876.9
> SS_E <- sum((n - 1) * vars)
> SS_E
[1] 725.2
> df_trt <- 2
> df_E   <- 12
> MS_trt <- SS_trt / df_trt
> MS_E   <- SS_E   / df_E
> F_stat <- MS_trt / MS_E
> F_stat
[1] 7.255102
> p_val  <- 1 - pf(F_stat, df_trt, df_E)
> p_val
[1] 0.00860219
> # 1(c): error SD
> sigma_hat <- sqrt(MS_E)
> sigma_hat
[1] 7.773888
> # 1(d): pairwise differences, Tukey HSD and Scheffé
> SE_diff <- sqrt(2 * MS_E / n)
> SE_diff
[1] 4.916638
> diffs <- c(
+   Priv_Pub = means["Private"] - means["Public"],
+   Ch_Pub   = means["Church"]  - means["Public"],
+   Ch_Priv  = means["Church"]  - means["Private"]
+ )
> diffs
Priv_Pub.Private    Ch_Pub.Church   Ch_Priv.Church
            17.4              2.7            -14.7
> qcrit <- qtukey(0.95, nmeans = 3, df = df_E)
> qcrit
[1] 3.772929
> HSD   <- qcrit * sqrt(MS_E / n)
> HSD
[1] 13.11692
> Fcrit <- qf(0.95, df1 = df_trt, df2 = df_E)
> Fcrit
[1] 3.885294
> k_S   <- (3 - 1) * Fcrit
> k_S
[1] 7.770588
> T2 <- (diffs^2) / (SE_diff^2)
> T2
Priv_Pub.Private    Ch_Pub.Church   Ch_Priv.Church
       12.524545         0.301572         8.939189
> ########################
> # Problem 2
> ########################
>
> # 2(a): One-way ANOVA from summary stats
> means_P2 <- c(G1 = 72.3, G2 = 79.5, G3 = 65.8)
> n_P2     <- 6
> a_P2     <- 3
> N_P2     <- a_P2 * n_P2
> overall_var <- 84.03
> MS_E_P2     <- 57.84
> grand_mean_P2 <- mean(means_P2)
> grand_mean_P2
[1] 72.53333
> SS_T_P2 <- overall_var * (N_P2 - 1)
> SS_T_P2
[1] 1428.51
> SS_trt_P2 <- n_P2 * sum((means_P2 - grand_mean_P2)^2)
> SS_trt_P2
[1] 563.56
> SS_E_P2_alt <- SS_T_P2 - SS_trt_P2
> SS_E_P2_alt
[1] 864.95
> df_trt_P2 <- a_P2 - 1
> df_E_P2   <- N_P2 - a_P2
> SSE_from_MS <- MS_E_P2 * df_E_P2
> SSE_from_MS
[1] 867.6
> SS_E_P2 <- SSE_from_MS
> SS_E_P2
[1] 867.6
> SS_T_P2_consistent <- SS_trt_P2 + SS_E_P2
> SS_T_P2_consistent
[1] 1431.16
> MS_trt_P2 <- SS_trt_P2 / df_trt_P2
> MS_trt_P2
[1] 281.78
> F_stat_P2 <- MS_trt_P2 / MS_E_P2
> F_stat_P2
[1] 4.871715
> p_val_P2 <- 1 - pf(F_stat_P2, df_trt_P2, df_E_P2)
> p_val_P2
[1] 0.023428
> # 2(c): pairwise differences vs control (group 2)
> SE_diff_P2 <- sqrt(2 * MS_E_P2 / n_P2)
> SE_diff_P2
[1] 4.3909
> diffs_vs_control <- c(
+   G1_minus_G2 = means_P2["G1"] - means_P2["G2"],
+   G3_minus_G2 = means_P2["G3"] - means_P2["G2"]
+ )
> diffs_vs_control
G1_minus_G2.G1 G3_minus_G2.G3
          -7.2          -13.7
> t_stats_vs_control <- diffs_vs_control / SE_diff_P2
> t_stats_vs_control
G1_minus_G2.G1 G3_minus_G2.G3
     -1.639755      -3.120090
> # 2(d): contrast = average of groups 1 and 3 vs control
> C_hat <- 0.5 * means_P2["G1"] + 0.5 * means_P2["G3"] - means_P2["G2"]
> C_hat
    G1
-10.45
> SE_C <- sqrt(MS_E_P2 * 1.5 / n_P2)
> SE_C
[1] 3.802631
> t_C <- C_hat / SE_C
> t_C
       G1
-2.748098
> alpha      <- 0.05
> K          <- 10
> alpha_star <- alpha / K
> alpha_star
[1] 0.005
> t_crit_Bonf <- qt(1 - alpha_star / 2, df = df_E_P2)
> t_crit_Bonf
[1] 3.286039
> ########################
> # Problem 3
> ########################
>
> # 3(a): One-way ANOVA on raw gene expression
> dat <- read.csv("Stats 490/Exam1/Mid1Prob3.csv")
> str(dat)
'data.frame':   30 obs. of  4 variables:
 $ X              : int  1 2 3 4 5 6 7 8 9 10 ...
 $ Project        : int  1 2 3 4 5 6 7 8 9 10 ...
 $ Treatment      : chr  "MP Only" "MP Only" "MP Only" "MP Only" ...
 $ Gene.Expression: num  334.5 31.6 701 41.2 61.2 ...
> head(dat)
  X Project Treatment Gene.Expression
1 1       1   MP Only           334.5
2 2       2   MP Only            31.6
3 3       3   MP Only           701.0
4 4       4   MP Only            41.2
5 5       5   MP Only            61.2
6 6       6   MP Only            69.6
> dat$Treatment <- factor(dat$Treatment)
> library(dplyr)
> group_stats <- dat %>%
+   group_by(Treatment) %>%
+   summarise(
+     n    = n(),
+     mean = mean(Gene.Expression),
+     sd   = sd(Gene.Expression)
+   )
> group_stats
# A tibble: 3 × 4
  Treatment      n  mean    sd
  <fct>      <int> <dbl> <dbl>
1 MP + HDMTX    10  479.  371.
2 MP + LDMTX    10  165.  126.
3 MP Only       10  238.  308.
> fit_raw <- aov(Gene.Expression ~ Treatment, data = dat)
> summary(fit_raw)
            Df  Sum Sq Mean Sq F value Pr(>F)
Treatment    2  538442  269221   3.249 0.0544 .
Residuals   27 2236974   82851
---
Signif. codes:  0 '***' 0.001 '**' 0.01 '*' 0.05 '.' 0.1 ' ' 1
> anova_raw <- summary(fit_raw)[[1]]
> anova_raw
            Df  Sum Sq Mean Sq F value  Pr(>F)
Treatment    2  538442  269221  3.2495 0.05439 .
Residuals   27 2236974   82851
---
Signif. codes:  0 '***' 0.001 '**' 0.01 '*' 0.05 '.' 0.1 ' ' 1
> # 3(b): normality check for raw residuals
> res_raw <- residuals(fit_raw)
> par(mfrow = c(1, 2))
> qqnorm(res_raw, main = "QQ Plot - Raw Data Residuals")
> qqline(res_raw)
> hist(res_raw, main = "Histogram - Raw Data Residuals",
+      xlab = "Residuals")
> par(mfrow = c(1, 1))
> shapiro_raw <- shapiro.test(res_raw)
> shapiro_raw

        Shapiro-Wilk normality test

data:  res_raw
W = 0.91654, p-value = 0.02183

> # 3(c): log transform and ANOVA
> dat$logGE <- log(dat$Gene.Expression)
> group_stats_log <- dat %>%
+   group_by(Treatment) %>%
+   summarise(
+     n    = n(),
+     mean = mean(logGE),
+     sd   = sd(logGE)
+   )
> group_stats_log
# A tibble: 3 × 4
  Treatment      n  mean    sd
  <fct>      <int> <dbl> <dbl>
1 MP + HDMTX    10  5.73  1.13
2 MP + LDMTX    10  4.74  1.00
3 MP Only       10  4.79  1.18
> fit_log <- aov(logGE ~ Treatment, data = dat)
> summary(fit_log)
            Df Sum Sq Mean Sq F value Pr(>F)
Treatment    2   6.30   3.148   2.571 0.0951 .
Residuals   27  33.07   1.225
---
Signif. codes:  0 '***' 0.001 '**' 0.01 '*' 0.05 '.' 0.1 ' ' 1
> anova_log <- summary(fit_log)[[1]]
> anova_log
            Df Sum Sq Mean Sq F value  Pr(>F)
Treatment    2  6.297  3.1485  2.5707 0.09507 .
Residuals   27 33.069  1.2248
---
Signif. codes:  0 '***' 0.001 '**' 0.01 '*' 0.05 '.' 0.1 ' ' 1
> # 3(d): residual diagnostics for log model
> res_log    <- residuals(fit_log)
> fitted_log <- fitted(fit_log)
> par(mfrow = c(1, 2))
> plot(fitted_log, res_log,
+      main = "Residuals vs Fitted (log model)",
+      xlab = "Fitted values", ylab = "Residuals")
> abline(h = 0, col = "red")
> qqnorm(res_log, main = "QQ Plot - log model residuals")
> qqline(res_log, col = "red")
> par(mfrow = c(1, 1))
> shapiro_log <- shapiro.test(res_log)
> shapiro_log

        Shapiro-Wilk normality test

data:  res_log
W = 0.95648, p-value = 0.2511
\end{verbatim}


\end{document}